
% =====================================================================
\section{Is the singlet superfluid the Higgs field?}
\label{sec:higgs}
% =====================================================================

Hypothesis H2 states that the superfluid singlet state of Sec.~\ref{sec:order}, the condensate of
the pair operator along the hidden axis, \emph{is the Higgs field}. This section meets objective O5.
It lists what the Higgs field of the Standard Model is required to do, tests the condensate against
each requirement, and states the result.

\subsection{What the Higgs field does, and the resulting conditions}

In the Standard Model the Higgs field $H$ is a doublet of complex scalar fields under the
weak-isospin group $SU(2)_L$, with weak hypercharge $\tfrac12$
\cite{Higgs1964,EnglertBrout1964,Weinberg1967}. Its potential is minimal at
$|H|=v_{\rm EW}/\sqrt2$ with $v_{\rm EW}=(\sqrt2G_F)^{-1/2}=246.22$~GeV. This vacuum value breaks
$SU(2)_L\times U(1)_Y$ to the $U(1)$ of electromagnetism, gives the $W$ and $Z$ bosons their masses
with the ratio $\rho\equiv M_W^2/(M_Z^2\cos^2\theta_W)=1$, and gives the quarks and charged leptons
their masses through Yukawa couplings. The one remaining scalar particle has the mass 125.20~GeV and
couplings proportional to the masses of the particles to which it couples \cite{PDG2024}.

Two structural facts are essential for what follows. The first is the \emph{custodial symmetry}. The
four real components of $H$ form a vector of $SO(4)$, which is locally $SU(2)_L\times SU(2)_R$, and
the Higgs potential depends only on its length. The vacuum value points in one direction and breaks
$SO(4)$ to $SO(3)$; that residual $SO(3)$ is what protects $\rho=1$ \cite{Sikivie1980}. In the
language of representations, $H$ is a \emph{real} $(2,2)$. The second fact is that the idea of a
composite Higgs field is old: it is the idea of Nambu and Jona-Lasinio \cite{NJL1961} applied to the
electroweak scale, with a condensate of top quarks \cite{BardeenHillLindner1990}. The present
hypothesis is of this kind, with one difference: the pair is made of two fermions and carries fermion
number 2, where those condensates are made of a fermion and an antifermion.

A condensate that is to \emph{be} the Higgs field must satisfy the seven conditions of
Table~\ref{tab:conditions}.

\begin{table}
\caption{\label{tab:conditions}Conditions that a condensate must satisfy in order to be the Higgs
field of the Standard Model.}
\begin{ruledtabular}
\begin{tabular}{l>{\raggedright\arraybackslash}p{0.47\columnwidth}>{\raggedright\arraybackslash}p{0.36\columnwidth}}
 & Condition & Reason \\
\hline
C1 & scalar under the Lorentz group of ordinary spacetime & the vacuum is Lorentz invariant \\
C2 & non-zero in empty space & the vacuum contains no medium \\
C3 & a doublet of $SU(2)_L$ with a neutral component & to break the electroweak symmetry and leave electromagnetism \\
C4 & symmetry breaking $SO(4)\to SO(3)$ & $\rho=1$ \\
C5 & scale 246 GeV; a scalar particle of 125 GeV with the observed couplings & measured \\
C6 & couplings that give the quarks and leptons their masses & measured \\
C7 & no further light or charged particles that would have been seen & measured \\
\end{tabular}
\end{ruledtabular}
\end{table}

\subsection{Condition C1: a scalar in ordinary spacetime}

This condition is met. The pair operator $\Delta^0=\psi^T\Cm\gamma^0\psi$ is the component along $x0$
of a vector of the eight-dimensional spacetime, and it is invariant under all Lorentz transformations
that do not involve $x0$ (Sec.~\ref{sec:hiddenpair}), hence under the Lorentz group of ordinary
spacetime. This is the property for which the hidden axis is needed: a scalar pair of the full
eight-dimensional group does not exist (Theorem \SIG{3}{4}), but a vector along a dimension that
ordinary physics does not see is a scalar for ordinary physics.

\subsection{Condition C2: a condensate in empty space}

This condition is not established. The superfluid forms in a \emph{medium} of fermions, and the Higgs
field must condense in the vacuum. The universe cannot be filled with the medium: a medium defines a
rest frame and would violate Lorentz invariance, and a Fermi sea with the energy scale of the weak
interaction would have an energy density of the order of $(100~\mathrm{GeV})^4$, some fifty-five
orders of magnitude above the density of the universe.

The pair can condense without a medium only if the attraction exceeds a critical strength, as in the
model of Nambu and Jona-Lasinio. Section~\ref{sec:vacuum} derived the condition in the ladder
approximation: $\aVS>\alpha_c$, with $\alpha_c={{GENM:cjt_alpha_c_10}}$ for a cutoff of ten fermion
masses, ${{GENM:cjt_alpha_c_100}}$ for a hundred and ${{GENM:cjt_alpha_c_1e6}}$ for $10^6$ (at
$m_s=m$), and $\alpha_c\to\pi/2$ for an infinite cutoff. A vacuum condensate therefore needs
$\aVS$ of about 2 or more, that is, $\gVS\gtrsim5$. At such couplings two fermions already form a
bound state, the potential of the Vorton-Scalar is unstable at one loop unless its self-coupling is
equally strong, and neither the one-boson-exchange kernel nor the ladder approximation is
quantitatively reliable. In the weak-coupling range, where the model has been solved, there is no
vacuum condensate.

\subsection{Conditions C3 and C4: the electroweak quantum numbers}

For a condensate of the field $\psi$ to break the electroweak symmetry, $\psi$ must carry electroweak
charges. The only symmetries of the model that commute with the Lorentz group of ordinary spacetime,
and that could therefore be internal symmetries of four-dimensional physics, are the symmetries of
Sec.~\ref{sec:internal}. The question is whether the electroweak group fits into them, with the
condensate in the right representation.

%%<1671>
On the $7+1$ spacetime the four extra axes are space-like, and their rotation group $SO(4)$, locally
$SU(2)_A\times SU(2)_B$, commutes with the ordinary rotations. For fields that do not depend on the
extra coordinates it is the complete non-Abelian internal symmetry (Table~\ref{tab:symmetry}). The
representations are those found in Secs.~\ref{sec:pairsrest} and \ref{sec:order}. One particle at
rest is $(2;2,1)\oplus(2;1,2)$ under $SO(3)\times SU(2)_A\times SU(2)_B$: it has spin $\tfrac12$ and
is a doublet of one of the two groups. The ten $s$-wave pairs of ordinary spin 0 are
$(2,2)\oplus(3,1)\oplus(1,3)$, and the $(2,2)$ are the four vector pairs $\Delta^a$, a four-vector of
$SO(4)$. The neutral current of the Vorton-Scalar attracts in the $(2,2)$ and repels in the
$(3,1)\oplus(1,3)$, and the state of lowest energy is $\Delta^a=|\Delta|e^{i\varphi}n^a$ with a real
unit vector $n^a$.

Compare this with the Higgs field. \emph{The order parameter is a $(2,2)$ of
$SU(2)_A\times SU(2)_B$, and its vacuum value breaks $SO(4)\to SO(3)$: the representation and the
breaking pattern of the Higgs field and of its custodial symmetry.} If $SU(2)_A$ is identified with
$SU(2)_L$ and the hypercharge with the generator $T^3_B$ of $SU(2)_B$, the four components of
$\Delta^a$ have the electric charges $Q=T^3_A+T^3_B=(+1,0,0,-1)$, the direction of the hidden axis is
electrically neutral, and conditions C3 and C4 are met. Four qualifications apply.

(i) \emph{The order parameter is complex.} $\Delta^a$ has four complex components, twice the content
of the Higgs doublet, because the pair carries fermion number. Its vacuum value also breaks $U(1)_F$.
If the fermion number is an exact global symmetry there is an additional massless Goldstone boson,
and besides the Higgs particle a second neutral and a charged scalar: the structure is that of a model
with two Higgs doublets.

(ii) \emph{The identification of the gauge group is a hypothesis.} Nothing in the model makes
$SU(2)_A$ a \emph{gauge} symmetry. In theories of the Kaluza-Klein type, rotations of extra
dimensions become gauge symmetries when the extra space is curled up with the corresponding isometry
group \cite{Witten1981}; a four-dimensional torus has only the isometries $U(1)^4$, a four-sphere
has $SO(5)\supset SO(4)$. Whether the second stage of the program, in which the hidden axis curls up
into a ring and the other three extra dimensions deflate, leaves an $SU(2)$ is open. A ring by itself
has only $U(1)$, and treating $x0$ differently from $x5,x6,x7$ breaks the $SO(4)$ that the custodial
structure needs.

(iii) \emph{In the vacuum the multiplet is larger still.} A chemical potential singles out the pairs
of fermion number 2. Without it the zero-mode sector has the symmetry $O(4)\times O(4)$, and the pair
condensate is one orientation of a 16-component real multiplet that also contains condensates of a
particle and an antiparticle (Sec.~\ref{sec:vacuum}). A vacuum condensate would have to be analyzed
in that larger space.

(iv) \emph{Ordinary matter must carry the same charges.} Quarks and leptons are doublets of $SU(2)_L$.
If $SU(2)_L$ is a rotation group of the extra dimensions, they must be spinors of those rotations,
that is, components of fields like $\psi$ itself. A picture in which the Standard Model, with an
elementary Higgs doublet, lives on a brane and has its own gauge groups is an alternative to the
present hypothesis, not a part of it.

On the $4+4$ spacetime of the companion paper \cite{Companion1644} the same algebraic pattern exists
for the zero modes, but there the group $SO(4)$ is a hidden symmetry that is not a transformation of
the spacetime, and the geometry cannot turn it into a gauge symmetry. The metric on the four extra
axes is $\mathrm{diag}(-,-,-,-)$ here and $\mathrm{diag}(+,-,-,-)$ there, and the compact group
$SO(4)$ is the isometry group only of the first.
%%</1671>
%%<1644>
On the $4+4$ spacetime the hidden axis is space-like and the three other extra axes are time-like,
and two symmetry groups must be distinguished.

\emph{The symmetry that the geometry provides.} The rotations of the three additional times,
$SO(3)_T$, are a symmetry of the quantized model whether or not the fields depend on $x0$. Under the
corresponding $SU(2)_T$ one particle at rest is a doublet, pairs are singlets or triplets, the
hidden-axis pair $\Delta^0$ is a singlet, and $\Delta^5,\Delta^6,\Delta^7$ form a triplet (checks
\texttt{44-Q1} to \texttt{44-Q3}). If the weak isospin were $SU(2)_T$, the hidden-axis condensate
would not break the electroweak symmetry at all, and a condensate of the triplet would break it with
$\rho\neq1$ at tree level, $\rho=\tfrac12$ for hypercharge 1 and no $Z$ mass at all for hypercharge
0, against the measured value 1 to a part in a thousand \cite{PDG2024}.

\emph{The hidden symmetry of the zero-mode sector.} For fields that do not depend on $x0$ either, the
model has the larger symmetry $SU(2)_A\times SU(2)_B$ of Sec.~\ref{sec:internal}, of which $SU(2)_T$
is the diagonal subgroup. The particles are $(2,1)\oplus(1,2)$, the four pairs
$(\Delta^0,i\Delta^5,i\Delta^6,i\Delta^7)$ form the $(2,2)$, and the condensate breaks
$SO(4)\to SO(3)_T$: algebraically the representation and the breaking pattern of the Higgs field and
of its custodial symmetry. But the three generators beyond $SO(3)_T$ are the matrices
$\gamma^0\gamma^t$, which act on the components of the field only. They are not rotations of the
spacetime; the transformations of the spacetime that mix $x0$ with a time are boosts, generated by the
same matrices without the factor $i$, and boosts are not unitary. Hence the hidden symmetry is not an
isometry of any extra space, so that the mechanism of Kaluza and Klein \cite{Witten1981} cannot make
it a gauge symmetry; it is violated by every mode of the fields that depends on $x0$; and it could be
gauged only by hand, in the four-dimensional theory of the zero modes, which would be a new postulate
unrelated to the geometry of the $4+4$ spacetime.

\emph{On the $4+4$ spacetime the algebraic pattern of the Higgs vacuum is present in the zero-mode
sector, but the electroweak gauge group has no geometric origin there.} This is the precise form of
the difference between the two spacetimes. The metric on the four extra axes is
$\mathrm{diag}(+,-,-,-)$ here and $\mathrm{diag}(-,-,-,-)$ on the $7+1$ spacetime of the companion
paper \cite{Companion1671}; the compact group $SO(4)$ is the isometry group only of the second, and
here it survives as a symmetry of the zero modes alone.
%%</1644>

\subsection{Conditions C5, C6 and C7: scale, masses and spectrum}

These conditions concern dynamics that the present model does not determine. What can be said is the
following.

\emph{C5, the scale and the mass of the scalar particle.} In models of the Nambu-Jona-Lasinio type
the scalar bound state that accompanies the condensate has a mass of about twice the mass that the
condensate gives to its constituent fermions. With 125 GeV for the scalar, the constituents would
weigh about 60 GeV. By C3 they carry electroweak charges, and particles of that mass with electroweak
charges would have been produced in pairs at the electron-positron collider LEP, whose searches
exclude new charged fermions below about 100 GeV, and in decays of the $Z$ boson if lighter than
45 GeV \cite{PDG2024}. Refinements of the relation between the two masses
\cite{BardeenHillLindner1990} change it logarithmically, not by the factor needed. This estimate
speaks against the hypothesis. It is not a proof, because the strong-coupling dynamics of C2 is not
under control.

\emph{C6, the masses of quarks and leptons.} The model contains no coupling of the pair field to
ordinary fermions. A pair of fermion number 2 cannot couple to a fermion-antifermion bilinear of
ordinary matter without violating either the fermion number or the conservation of the ordinary
quantum numbers; an additional mechanism would have to be invented. This condition is not addressed.

\emph{C7, additional particles.} \SIG{If the hidden symmetry were gauged by hand, the model would
predict}{The model predicts}, besides the Higgs particle, the partners that complete the complex
$(2,2)$, a second neutral and a charged scalar; a Goldstone boson of the fermion number; the fermions
$\psi$ themselves, with electroweak charges; and the two modes of the Vorton-Scalar. None of these is
observed. This is an obstacle, of a size that depends on unknown masses.

\emph{The three hypotheses need three different scales.} The Higgs hypothesis needs the condensate at
246 GeV. The analysis of dark matter in Sec.~\ref{sec:dm} points to bound pairs heavier than about
0.2 GeV with a coupling of nuclear strength, or to no superfluid at all. The mechanism for the Hubble
tension of Sec.~\ref{sec:hubble} needs fermions of about 1 eV. One field with one mass cannot serve
all three purposes.

\subsection{Verdict}

\begin{table*}
\caption{\label{tab:verdict}The Higgs hypothesis H2 against the seven conditions of
Table~\ref{tab:conditions}, for the spacetime of this paper and for that of the companion paper.}
\begin{ruledtabular}
\begin{tabular}{>{\raggedright\arraybackslash}p{0.21\textwidth}>{\raggedright\arraybackslash}p{0.24\textwidth}>{\raggedright\arraybackslash}p{0.27\textwidth}>{\raggedright\arraybackslash}p{0.2\textwidth}}
Condition & $7+1$ & $4+4$ (quantized sector) & Basis \\
\hline
C1, scalar in ordinary spacetime & met & met & Sec.~\ref{sec:hiddenpair} \\
C2, condensate in the vacuum & open; needs $\aVS\gtrsim2$ & open; needs $\aVS\gtrsim2$ & ladder equation, Sec.~\ref{sec:vacuum} \\
C3, electroweak doublet with neutral component & met, \emph{if} $SU(2)_L=SU(2)_A$ and $Y=T^3_B$ & only under a hidden symmetry that the geometry cannot gauge; singlet under $SO(3)_T$ & Secs.~\ref{sec:pairsrest}, \ref{sec:internal} \\
C4, breaking $SO(4)\to SO(3)$ & met & met in the zero-mode sector only & Secs.~\ref{sec:internal}, \ref{sec:ground} \\
C5, scale and scalar mass & unfavorable (estimate) & --- & this section \\
C6, masses of ordinary fermions & not addressed & --- & this section \\
C7, no unobserved particles & additional particles predicted & --- & this section \\
\end{tabular}
\end{ruledtabular}
\end{table*}

Table~\ref{tab:verdict} collects the results.
%%<1671>
The hypothesis ``the singlet superfluid is the Higgs field'' is \emph{unproven on the $7+1$
spacetime, with one structural point in its favor and two estimates against it.} The structural
point is not trivial: the pair is a scalar of ordinary spacetime, it belongs to a $(2,2)$ of the
rotation group of the extra space axes, and its condensate has the custodial breaking pattern of the
Higgs vacuum. Against it stand the estimate of the scale and of the constituent mass (C5) and the
predicted additional particles (C7), and the fact that a condensate in the vacuum has been found only
in a strong-coupling regime that the present methods do not control (C2).
%%</1671>
%%<1644>
The hypothesis ``the singlet superfluid is the Higgs field'' is \emph{without a geometric basis on the
$4+4$ spacetime.} The pair is a scalar of ordinary spacetime, and the algebraic pattern of the Higgs
vacuum exists for the zero modes, but only under a hidden symmetry. Under the rotations of the
additional times, the only non-Abelian internal symmetry that the geometry of $4+4$ provides, the
condensate is a singlet and cannot break the electroweak symmetry. On the $7+1$ spacetime of the
companion paper the same pattern is geometric, and the hypothesis is unproven there, with one
structural point in its favor and two estimates against it \cite{Companion1671}.
%%</1644>

A remark on what the model \emph{does} contain: the Vorton-Scalar is itself a scalar field with a
vacuum value that gives the fermions $\psi$ their mass. It plays for these fermions the role that the
Higgs field plays for ordinary matter. But it is a singlet of all non-Abelian symmetries, and it
cannot be the Higgs field of the Standard Model either.

\subsection{What has been established, and what has not}

Established: the hidden-axis pair is a scalar of ordinary spacetime.
\SIG{For the zero modes it belongs to a $(2,2)$ of a hidden $SU(2)_A\times SU(2)_B$, and its
condensate breaks $SO(4)\to SO(3)_T$; under the rotations of the additional times it is a
singlet.}{It belongs to a $(2,2)$ of the rotation group of the four extra space axes, and its
condensate has the custodial breaking pattern $SO(4)\to SO(3)$ of the Higgs vacuum.}

Not established: that a condensate forms in the vacuum; that
\SIG{any symmetry of the model}{the rotation group of the extra dimensions} is the electroweak gauge
group; that the scale, the mass of the scalar particle and its couplings come out as observed; and
that ordinary fermions obtain their masses.
